\chapter{Ανάλυση Σημαντικότητας Χαρακτηριστικών}
\label{ch:fi}

Η ανάλυση σημαντικότητας χαρακτηριστικών (feature importance analysis) προσφέρει
ερμηνευσιμότητα στα μοντέλα ML, αποκαλύπτοντας ποιοι παράγοντες οδηγούν τις
προβλέψεις. Στην παρούσα εργασία, εξετάζουμε πώς η στρατηγική πρόβλεψης
(TF/Recursive/MIMO/Direct) αλλάζει τη βαρύτητα που αποδίδουν τα μοντέλα σε
χρονικές υστερήσεις, εξωγενείς παραμέτρους και ημερολογιακά χαρακτηριστικά.

% ─────────────────────────────────────────────────────────────────────────────
\section{Μεθοδολογία}
\label{sec:fi_methodology}
% ─────────────────────────────────────────────────────────────────────────────

\subsection*{Gain-Based Feature Importance (LGBM / XGBoost)}

Για τους αλγόριθμους LightGBM και XGBoost, χρησιμοποιείται η σημαντικότητα
βάσει \textit{κέρδους πληροφορίας} (gain): σε κάθε διαχωριστικό κόμβο δένδρου,
το gain μετρά τη μέση μείωση της αντικειμενικής συνάρτησης (MSE) λόγω του
συγκεκριμένου χαρακτηριστικού, σταθμισμένη από τον αριθμό δειγμάτων που
περνούν από τον κόμβο. Η κανονικοποιημένη αθροιστική συνεισφορά κάθε
χαρακτηριστικού προκύπτει ως~\cite{ke2017lgbm}:
\begin{equation}
    \text{Importance}(f) = \frac{\sum_{t} \text{gain}(t,f) \cdot \mathbb{1}[\text{feature}(t)=f]}
    {\sum_{f'} \sum_{t} \text{gain}(t,f')}
\end{equation}

\subsection*{Gini Impurity Importance (Random Forest)}

Για το Random Forest, χρησιμοποιείται η μέση μείωση ακαθαρσίας Gini (MDI) σε
όλα τα δένδρα του δάσους~\cite{breiman2001rf}. Η μέθοδος αυτή τείνει να
ευνοεί χαρακτηριστικά με υψηλή πληροφορία (όπως συνεχείς μεταβλητές με πολλές
τιμές), γεγονός που λαμβάνεται υπόψη στην ερμηνεία.

\subsection*{Εξαγωγή για 15 Παραλλαγές}

Εξάγεται feature importance για 15 μοντέλα, κατηγοριοποιημένα ανά στρατηγική:
\begin{itemize}
    \item \textbf{TF/CL:} LGBM-CL, XGB-CL, RF-CL (152 χαρακτηριστικά τιμής,
    133 φορτίου)
    \item \textbf{Recursive/OL:} LGBM-OL-Daily-Optuna, XGB-OL-SS-Dense
    \item \textbf{MIMO Direct:} LGBM-Direct-Dense (170 χαρακτηριστικά)
    \item \textbf{MIMO Multi-output:} XGB-MIMO-Dense, RF-MIMO-two-stage
\end{itemize}

% ─────────────────────────────────────────────────────────────────────────────
\section{Αποτελέσματα — Τιμή DAM}
\label{sec:fi_price}
% ─────────────────────────────────────────────────────────────────────────────

\begin{figure}[htbp]
    \centering
    \includegraphics[width=\textwidth]{ch4_feature_importance/fi_viz_price.png}
    \caption{Top-15 χαρακτηριστικά για τα μοντέλα τιμής ανά αλγόριθμο και
    στρατηγική. Η αδρή στήλη αντιστοιχεί στη σημαντικότητα gain/MDI
    (κανονικοποιημένη στο 100\%).}
    \label{fig:fi_price}
\end{figure}

\subsection*{Teacher-Forcing / Closed-Loop}

Στα μοντέλα TF/CL, παρατηρείται έντονη \textit{συγκέντρωση} στη χρονική
υστέρηση $y_{t-1}$:
\begin{itemize}
    \item \textbf{LGBM-CL}: $y_{t-1} = 88,3\%$ της συνολικής σημαντικότητας.
    Ακολουθούν $y_{t-24}$ (3,1\%), gas\_price (2,0\%) και $y_{t-168}$ (1,4\%).
    \item \textbf{XGB-CL}: Λιγότερο ακραία συγκέντρωση· $y_{t-1} = 52\%$,
    $y_{t-24} = 12\%$, gas\_price = 8\%.
    \item \textbf{RF-CL}: Πιο ομοιόμορφη κατανομή· $y_{t-1} = 28\%$, με
    σημαντική παρουσία $y_{t-24}$, $y_{t-48}$, gas\_price.
\end{itemize}

Η κυριαρχία του $y_{t-1}$ στο LGBM-CL είναι αναμενόμενη: στη στρατηγική TF,
οι πραγματικές lagged τιμές παρέχονται στο μοντέλο, και η τιμή ηλεκτρισμού
παρουσιάζει σημαντική αυτοσυσχέτιση (ρ ≈ 0,87 για lag-1).

\subsection*{Recursive / Open-Loop}

Στα OL μοντέλα, η σημαντικότητα ανακατανέμεται:
\begin{itemize}
    \item \textbf{LGBM-Daily-Optuna}: $y_{t-1}$ = 41\%, $y_{t-24}$ = 18\%,
    $y_{t-168}$ = 12\%, gas\_price = 9\%. Η μειωμένη βαρύτητα του $y_{t-1}$
    σε σχέση με CL υποδηλώνει ότι το μοντέλο δεν εμπιστεύεται πλήρως την
    πιο πρόσφατη πρόβλεψη (η οποία περιέχει σφάλμα).
    \item \textbf{XGB-SS-Dense-OL}: Με τα dense lags (lags 4--23), οι ενδοημερήσιες
    τιμές αποκτούν συλλογική βαρύτητα ~25\%, αντικατοπτρίζοντας ότι η ενδοημερήσια
    δυναμική αξιοποιείται περισσότερο.
\end{itemize}

\subsection*{MIMO / Direct}

Στα MIMO/Direct μοντέλα, η δομή της σημαντικότητας αλλάζει ριζικά:
\begin{itemize}
    \item \textbf{LGBM-Direct-Dense}: Ανάδυση αιολικής (\textit{wind\_gen\_lag1}),
    ηλιακής (\textit{solar\_gen\_lag1}) και gas\_price ως top-4 χαρακτηριστικά.
    Η εξάρτηση από $y_{t-1}$ πέφτει στο 15--20\%, αφού το μοντέλο δεν έχει
    πρόσβαση σε ενδιάμεσες TF τιμές.
    \item \textbf{RF-MIMO-two-stage}: Κυριαρχεί το \textit{y\_roll24} (rolling
    mean 24h) με 51,6\%. Η κυριαρχία αυτή υποδηλώνει ότι το RF αξιοποιεί κυρίως
    το επίπεδο τιμής της προηγούμενης ημέρας για να «σταθεροποιήσει» τις
    προβλέψεις σε ένα εύλογο εύρος.
\end{itemize}

% ─────────────────────────────────────────────────────────────────────────────
\section{Αποτελέσματα — Ηλεκτρικό Φορτίο}
\label{sec:fi_load}
% ─────────────────────────────────────────────────────────────────────────────

\begin{figure}[htbp]
    \centering
    \includegraphics[width=\textwidth]{ch4_feature_importance/fi_viz_load.png}
    \caption{Top-15 χαρακτηριστικά για τα μοντέλα φορτίου ανά αλγόριθμο και
    στρατηγική.}
    \label{fig:fi_load}
\end{figure}

Το ηλεκτρικό φορτίο παρουσιάζει διαφορετικό προφίλ feature importance σε
σχέση με την τιμή:

\begin{itemize}
    \item \textbf{Μικρότερη συγκέντρωση} στο $y_{t-1}$: για LGBM-CL φορτίου,
    $y_{t-1} = 45\%$ (έναντι 88\% για τιμή), με σημαντική παρουσία
    $y_{t-24} = 20\%$, $y_{t-168} = 8\%$ και \textit{hour\_of\_day} = 6\%.
    Αυτό αντικατοπτρίζει τη σαφή ημερήσια εποχικότητα του φορτίου.

    \item \textbf{Ημερολογιακά χαρακτηριστικά} πιο σημαντικά: \textit{hour},
    \textit{day\_of\_week}, \textit{is\_holiday} εμφανίζονται σε top-10 για
    όλους τους αλγόριθμους φορτίου — αναμενόμενο λόγω εργασιακών κύκλων.

    \item \textbf{Θερμοκρασιακή επίδραση}: Λόγω απουσίας θερμοκρασιακών
    δεδομένων στο dataset, η έμμεση επίδρασή της αποτυπώνεται μέσω των
    rolling means ($y\_roll24$, $y\_roll168$) που λειτουργούν ως proxy.
\end{itemize}

% ─────────────────────────────────────────────────────────────────────────────
\section{Σύγκριση Κατά Στρατηγική}
\label{sec:fi_strategy}
% ─────────────────────────────────────────────────────────────────────────────

\begin{figure}[htbp]
    \centering
    \includegraphics[width=\textwidth]{ch4_feature_importance/fi_viz_strategy_shift.png}
    \caption{Μεταβολή σημαντικότητας τοπ-5 χαρακτηριστικών μεταξύ
    TF, OL και MIMO για τα μοντέλα τιμής (LGBM). Παρατηρείται
    σταδιακή μεταφορά βαρύτητας από $y_{t-1}$ σε $y_{t-24}$, gas\_price
    και αιολική παραγωγή καθώς εναλλάσσεται η στρατηγική.}
    \label{fig:fi_strategy}
\end{figure}

\Cref{fig:fi_strategy} αποτυπώνει τη διαδοχική μεταβολή της σημαντικότητας
χαρακτηριστικών. Τρεις κανονικότητες αναδύονται:

\begin{enumerate}
    \item \textbf{TF $\to$ Recursive}: Μείωση $y_{t-1}$ (88\% $\to$ 41\%), αύξηση
    $y_{t-24}$ και $y_{t-168}$. Το μοντέλο αντισταθμίζει την αβεβαιότητα
    της τελευταίας πρόβλεψης εστιάζοντας στις εποχικές τιμές.

    \item \textbf{Recursive $\to$ MIMO}: Περαιτέρω μείωση $y_{t-1}$ (41\% $\to$ 15--20\%),
    ανάδυση εξωγενών (gas\_price, αιολική). Χωρίς αναδρομικές υπολογισμένες
    τιμές ως είσοδο, το MIMO βασίζεται σε «εξωτερικές» πληροφορίες.

    \item \textbf{Dense Lags}: Τα intraday lags 4--23 εμφανίζονται συλλογικά ως
    δεύτερη ομάδα σημαντικότητας (~25\%) στα dense MIMO μοντέλα, αποδεικνύοντας
    την αξία της πλήρους ενδοημερήσιας ιστορίας.
\end{enumerate}

% ─────────────────────────────────────────────────────────────────────────────
\section{Κατηγοριοποίηση Χαρακτηριστικών}
\label{sec:fi_categories}
% ─────────────────────────────────────────────────────────────────────────────

\begin{figure}[htbp]
    \centering
    \includegraphics[width=\textwidth]{ch4_feature_importance/fi_viz_categories.png}
    \caption{Κατανομή σημαντικότητας ανά κατηγορία χαρακτηριστικών (price lags,
    calendar features, exogenous/market features, rolling statistics)
    για 8 παραλλαγές μοντέλων τιμής.}
    \label{fig:fi_categories}
\end{figure}

Τα χαρακτηριστικά κατηγοριοποιούνται σε τέσσερις ομάδες:

\begin{description}
    \item[Lagged values:] Χρονικές υστερήσεις $y_{t-k}$ (k=1,2,24,48,168 για
    standard· k=1--23 για dense). Κυριαρχούν στα TF/CL μοντέλα (>70\%).
    \item[Calendar features:] \textit{hour\_of\_day}, \textit{day\_of\_week},
    \textit{month}, \textit{is\_holiday}, \textit{is\_weekend}. Σχετικά σταθερή
    συνεισφορά (5--15\%) σε όλες τις στρατηγικές.
    \item[Market/exogenous:] gas\_price, CO\textsubscript{2}\_price,
    load\_fc (για τιμή), ανανεώσιμα lag1/lag24. Αυξάνει σε MIMO (15--25\%).
    \item[Rolling statistics:] \textit{y\_roll24}, \textit{y\_roll168}. Έμμεση
    αποτύπωση επιπέδου τιμής· σημαντικά για RF.
\end{description}

Κύριο εύρημα: καθώς μεταβαίνουμε από TF $\to$ Recursive $\to$ MIMO, η εξάρτηση από
lagged values μειώνεται (~70\% $\to$ 60\% $\to$ 40\%), ενώ η βαρύτητα market/exogenous
και calendar αυξάνεται αντίστοιχα. Αυτό υποδηλώνει ότι στην MIMO στρατηγική τα
μοντέλα αξιοποιούν σε μεγαλύτερο βαθμό τη δομή της αγοράς και τους εξωγενείς
παράγοντες, έναντι της αυτόματης αναδρομής.

% ─────────────────────────────────────────────────────────────────────────────
\section{Συμπεράσματα Κεφαλαίου}
\label{sec:fi_conclusions}
% ─────────────────────────────────────────────────────────────────────────────

Η ανάλυση feature importance αποκαλύπτει τρεις βασικές κανονικότητες:

\begin{enumerate}
    \item \textbf{Η στρατηγική καθορίζει τη σημαντικότητα}: Το ίδιο χαρακτηριστικό
    ($y_{t-1}$) έχει ριζικά διαφορετική βαρύτητα ανάλογα με τη στρατηγική (88\%
    στο TF vs 15\% στο MIMO), γεγονός που υποδηλώνει ότι τα μοντέλα «μαθαίνουν»
    να αξιοποιούν διαφορετικές ιδιότητες των δεδομένων.

    \item \textbf{Dense lags ανεβάζουν τη βαρύτητα ενδοημερήσιων χαρακτηριστικών}:
    Τα lags 4--23 προσθέτουν μια νέα κατηγορία πληροφορίας που ο αλγόριθμος
    αξιοποιεί σημαντικά, εξηγώντας τη βελτίωση απόδοσης MIMO.

    \item \textbf{Φορτίο vs τιμή}: Το φορτίο εμφανίζει πιο ισόρροπη κατανομή
    σημαντικότητας με ισχυρότερη παρουσία ημερολογιακών χαρακτηριστικών, αντανακλώντας
    την πιο ντετερμινιστική εποχική συμπεριφορά του έναντι της τιμής.
\end{enumerate}
